\documentclass[11pt,a4paper]{article}
\usepackage[utf8]{inputenc}
\usepackage[T1]{fontenc}
\usepackage{amsfonts}
\usepackage[french]{babel}
\frenchbsetup{StandardLists=true} % à inclure si on utilise \usepackage[francais]{babel}
\usepackage{enumitem}
\usepackage{amssymb}
\def\nbR{\ensuremath{\mathrm{I\! R}}}
\usepackage{amsmath}
\DeclareMathOperator{\e}{e}

\usepackage[left=2cm,right=2cm,top=2cm,bottom=2cm]{geometry}
\usepackage{multirow}
\usepackage[dvipsnames]{xcolor}
\usepackage{parskip}
\usepackage{caption}





\usepackage[T1]{fontenc}
\usepackage{fouriernc}

\usepackage{graphicx}
\usepackage{setspace}
\singlespacing

\usepackage{tcolorbox}
\usepackage{framed}
\usepackage{multicol}

\usepackage{fancybox}
\usepackage{fancyhdr}
\pagestyle{fancy}
\renewcommand\headrulewidth{1pt}
\fancyhead[L]{Lycée Denis Diderot }
\fancyhead[R]{Classe 3PM}
\setcounter{page}{1}\fancyfoot[C]{page \thepage}
\fancyfoot[R]{ \today}
\fancyfoot[L]{Alexis RUIZ}
\usepackage{pstricks,pst-plot}
\usepackage{tikz}
\usepackage{tkz-tab}
\usetikzlibrary{arrows}
\usepackage{pifont}

\newcommand{\ud}{\mathrm{d}}
\newcommand{\V}{\overrightarrow}
\newcommand{\Rep}{(O;\V{u};\V{v})} 
\newcommand{\cadreblanc}[1]{%
\medskip\framebox[\linewidth][c]{%
\rule{0mm}{#1mm}}\par
\medskip}
\setlength{\columnseprule}{.25mm}
\usepackage{multirow,tabularx,booktabs}
\usepackage{colortbl}
\usepackage{wrapfig}

\usepackage{qrcode}
\usepackage[tikz]{bclogo} 
\usepackage{ProfCollege}
\usepackage{pgfplots,mathtools}
\usepackage{bclogo}

\newcommand{\cadre}[1]
{\begin{bclogo}[couleur = white,logo=,  arrondi = 0.3,barre=none]
{%
\rule{0mm}{#1mm}}\par
\medskip
\end{bclogo}}



\begin{document}
\begin{center} 
\LARGE{C'est pas sorcier : l'aventure de l'électricité}  
\vspace{2mm}
\end{center}
\hrule
\normalsize
\vspace{2mm}
\begin{figure}[h]

{\begin{minipage}{0.59\textwidth}
NOM : \\

Prénom : 

\vspace{5mm}
\end{minipage}
}
{\begin{minipage}{0.4\textwidth}
\begin{center}
\qrcode[height=1in]{https://www.youtube.com/watch?v=efQW-ZmpyZs}
\end{center}
\end{minipage}}

\end{figure}


\vspace{2mm}
\hrule

\vspace{5mm}
Répondre aux questions ci-dessous en utilisant la vidéo en lien avec le QR Code : 
\begin{enumerate}
\item Qu’utilisait-on pour s’éclairer avant l’invention de l’électricité 
\cadre{15}
\item  Qu’a découvert Thalès de Milet ? 
\cadre{10}
\item Qu’est-ce-qu’un atome 
\cadre{15}
\item Qui appartient à la « famille plus » ? 
\cadre{15}
\item Que faut-il pour faire passer le courant ?
\cadre{15}
\item Qu’a inventé Benjamin Franklin ?
\cadre{15}
\item Qui a inventé la pile ?
\cadre{10} 
\newpage
\item De quoi une pile est-elle constituée ?
\cadre{15}
\item Quel est l’avantage des piles ? 
\cadre{15}
\item Quelles sont les deux catégories de courant ?
\cadre{15}
\item Que donnent des électrons qui bougent ? 
\cadre{15}
\item Comment s’appelle la roue avec des aimants qui fabrique de l’électricité ?
\cadre{15}
\item Quelles sont les différentes usines dans lesquelles on peut fabriquer de l’électricité ?
\cadre{15}
\item Qui a inventé l’ampoule ? 
\cadre{10}
\item Pourquoi y-a-t-il du verre autour d’une ampoule ?
\cadre{15}
\item Notre corps peut-il être « conducteur » de courant ?
\cadre{15}

\end{enumerate}


\end{document}
